\section{Methodology}
\label{sec:method}

\subsection{Problem Setup and Hypotheses}

A model is fine-tuned on tasks $T_1, \dots, T_6$ one after another, each for a \emph{stage} of 250 optimizer steps.
For a task learned at serial position $p$, we track its performance during all later stages.
The \emph{lag} is the amount of subsequent training, in steps or in stages.
We test three pre-registered hypotheses.

\para{H1 (shape).} Retention $R(t)$ after $t$ interfering steps follows a power law, not a linear (``uniform'') or exponential decay.
We count H1 as supported if the power law has the lowest AICc~\cite{hurvich1989regression} in a clear majority of task-level mean curves \emph{and} beats a two-parameter saturating exponential.

\para{H2 (domain).} Reasoning tasks are forgotten 3--5$\times$ faster than pattern-matching tasks.
We count H2 as supported if the reasoning/pattern ratio of mean forgetting has a bootstrap 95\% interval that excludes 1 and lies inside or overlaps $[3, 5]$, consistently across normalizations.

\para{H3 (position).} Early-learned tasks are forgotten disproportionately.
End-of-sequence retention is necessarily confounded with lag, because earlier tasks see more interference.
We count H3 as supported only if position has an effect on retention at \emph{fixed} lag.

\subsection{Models and Training}

Our main model is \llama (\texttt{meta-llama/Llama-3.2-1B-Instruct})~\cite{meta2024llama32}; we replicate on \qwen (\texttt{Qwen/Qwen2.5-1.5B-Instruct})~\cite{qwen2024qwen25}.
We use naive sequential fine-tuning with one LoRA adapter~\cite{hu2022lora} carried through all tasks ($r{=}16$, $\alpha{=}32$, no dropout, applied to all attention and MLP projections).
We train with AdamW~\cite{loshchilov2019decoupled} at learning rate $2\times10^{-4}$, no weight decay, batch size 16, and gradient clipping at 1.0.
The learning rate is constant after a 10-step linear warm-up at the start of every stage; we choose a constant rate so that a schedule does not modulate interference per step.
Each stage is 250 steps (2 epochs of 2{,}000 examples), so a run is 1{,}500 steps.
We use no replay and no regularization.
Decoding is greedy.

\subsection{Tasks}

\Tabref{tab:tasks} lists the six tasks, three per category.
The reasoning/pattern assignment is a judgment call, so we also report every result per task.

\begin{table}[t]
    \centering
    \small
    \resizebox{\textwidth}{!}{%
    \begin{tabular}{@{}llll@{}}
        \toprule
        \textbf{Task} & \textbf{Category} & \textbf{Source} & \textbf{Metric} \\
        \midrule
        \tgsm & reasoning & GSM8K~\cite{cobbe2021training}, worked solution + \texttt{Answer: n} & exact match \\
        \tarith & reasoning & synthetic 3-step integer arithmetic, worked steps & exact match \\
        \tlogic & reasoning & ProofWriter~\cite{tafjord2021proofwriter} (open-world, depth 1--3), 3-way & accuracy \\
        \tmt & pattern & WMT16 en$\rightarrow$de~\cite{bojar2016findings} & chrF~\cite{popovic2015chrf} \\
        \tsumm & pattern & SAMSum~\cite{gliwa2019samsum} & ROUGE-L F1~\cite{lin2004rouge} \\
        \temo & pattern & dair-ai/emotion~\cite{saravia2018carer}, 6 classes & accuracy \\
        \bottomrule
    \end{tabular}
    }
    \caption{The six tasks. Each has 2{,}000 training and 200 test examples. \temo gives the pattern category an accuracy-scored task, so category is not perfectly confounded with metric type.}
    \label{tab:tasks}
\end{table}


\subsection{Order Design and Evaluation}

\para{Orders.}
We use a 6-order Williams design~\cite{williams1949experimental}: a balanced Latin square in which every task occupies every position once and follows every other task once.
For \llama we run two blocks of six orders (seed 0 with the canonical square, seed 1 with a randomly relabelled square), giving 12 runs.
For \qwen we run one block of 6 runs.

\para{Evaluation.}
We measure held-out per-token loss on all six test sets every 25 steps and generation metrics every 125 steps, plus a zero-shot evaluation at step 0.
The shape analysis therefore rests on loss trajectories (51 points along the lag axis for a position-1 task), with the generation metrics as a coarser 11-point check.

\para{Baselines.}
We compare against (i) the zero-shot base model, (ii) each task's score at the end of its own stage, and (iii) joint multi-task training for 1{,}500 steps on the mixture of all six tasks (\llama: 2 seeds; \qwen: 1 seed).
Joint training tells us whether a shortfall of sequential training is forgetting or a capacity limit of the adapter.

\para{Onset experiment (exploratory).}
The 25-step grid turned out to be too coarse to resolve the onset of forgetting.
We therefore added a step-resolution experiment on \llama for six (old $\leftarrow$ new) task pairs.
We train the old task for 250 steps and then the new task for 100 steps, measuring old-task loss after every step up to step 20, every 2 steps up to step 50, and every 10 steps up to step 100.
This experiment uses one seed, and we chose the pairs partly because they showed strong interference in the main runs.

\subsection{Retention Measures}

Let $m_0$ and $L_0$ be a task's zero-shot metric and loss, $m^*$ and $L^*$ the values at the end of its own stage, and $m$ and $L$ the current values.
We use three pre-registered retention measures:
\begin{equation}
    \rloss = \frac{L_0 - L}{L_0 - L^*}, \qquad
    \rrel = \frac{m}{m^*}, \qquad
    \rgain = \frac{m - m_0}{m^* - m_0}.
\end{equation}
\rloss is our primary measure for shape because it is continuous and metric-free; it equals 1 when the task was just learned, 0 at zero-shot level, and can go negative.
\rgain is the fraction of learned gain that is kept; we compute it only where the gain $m^* - m_0$ is at least 0.05.
We also report the absolute drop $\absdrop = m^* - m$.
After seeing that a few cells far below zero-shot level dominate the unclipped mean, we added \rloss clipped to $[0, 1]$; we label it exploratory.
Forgetting is $1 - R$ (or \absdrop).

\subsection{Statistical Analysis}

\para{H1.}
We fit four curves to retention by least squares: linear $1 - at$, exponential $e^{-t/\tau}$, power $(1 + t/\tau)^{-b}$, and exponential-to-asymptote $c + (1-c)e^{-t/\tau}$.
The last is a fair two-parameter competitor to the power law.
We compare models by AICc on mean curves and on individual trajectories.

\para{H2.}
We compute the ratio of category means of forgetting over all stage-end evaluations with lag $\geq 1$, weighting tasks equally.
We obtain 95\% intervals from a cluster bootstrap over runs (10{,}000 resamples) and apply a Holm correction~\cite{holm1979simple} across the three pre-registered measures.
We also run an exact permutation test over the 20 ways of labelling 3 of the 6 tasks ``reasoning'', which treats tasks as the unit of generalization.

\para{H3.}
We compute the Spearman correlation between serial position and forgetting at lag 1 (after task-demeaning), and an OLS regression of forgetting on $\log(1 + \text{lag})$, position, and task with run-clustered errors.
In this parameterization a negative position coefficient means that earlier-learned tasks show more forgetting at the same lag, which is the direction H3 predicts.

We list all deviations from the pre-registered plan in \secref{app:deviations}.
